\section{Closed: two distinct primes}\label{sec:distinct}

The obvious strengthening of Conjecture~\ref{goldbach:conj} asks for the two
primes to be distinct.  It is false, and the counterexamples are the two
smallest cases.  It is recorded here, in the companion, rather than deleted:
a closed route keeps its exact failure mode searchable, and the corrected form
is what any later argument should quote.

\begin{prop}\label{distinct:counter}
There is no \term{goldbach-pair} $(p,q)$ with $p\neq q$ for $n=4$ or $n=6$.
For every even $n$ with $8\leq n\leq10^{6}$ such a pair does exist.
\end{prop}

\begin{proof}
For $n=4$ the only pair of primes summing to $4$ is $(2,2)$, since $3$ is prime
but $1$ is not.  For $n=6$ the candidates are $(2,4)$ and $(3,3)$, and $4$ is
not prime, so only $(3,3)$ survives.  Both fail the distinctness requirement.
The positive range is the output of
\begin{verbatim}
    python3 tools/python/distinct_goldbach.py 1000000
\end{verbatim}
which applies Lemma~\ref{reflect:lem} with the diagonal $p=n/2$ removed; the
retained output is \texttt{tools/logs/distinct-primes-1e6.log}.
\end{proof}

\emph{Why it fails, and what survives.}  The obstruction is not arithmetic but
trivial: the diagonal pair is available exactly when $n/2$ is prime, and for
$n\in\{4,6\}$ it is the only pair there is.  The statement is therefore true
for all $n\geq8$, and it is that corrected form, not the original, which should
be carried forward if the strengthening is ever needed.  The \term{sifted-count}
of Definition~\ref{sifted:def} is insensitive to the diagonal, so no
sieve-theoretic route would have detected this failure --- which is the reason
the check was run at all.
